% Nouveau EGA V (cote 25): archive cote 25, site page 55 (PDF page 56).
% typescript. Draft compared: none.
% AI-transcribed by Claude Opus 5.5 (claude-opus-5-5, adaptive thinking, effort high), Message Batches API, batch ega25b (msgbatch_01A7n5YuXubUPNK1zsepi8RG), 2026-11-13; not yet checked by eye.
\scanpage{25/55}
\mspage{50}

9. \uline{Changement d'immersion projective}.

9.1. Pour tout entier $n > 0$, soit $P(n) = P(\underline{Sym}^n(E))$, on a une
immersion évidente $u_n\colon P \to P(n)$, puisque $\underline{O}_P(n)$ est engendré par
ses sections au dessus de tout ouvert affine de S et que $p_*(\underline{O}_P(n))$
$\underline{Sym}^n(\underline{E})$, où $p\colon P \to S$ est la projection. S\var{\xout{o}}{}{typewriter strike-out}{ega25_25_055-v0001}i\var{\xout{t}}{}{typewriter strike-out}{ega25_25_055-v0002} $f\colon X\to P$ est un morphisme
non ramifié (resp. une immersion), il en est de même de $u_n f\colon X \to P(n)$.
Il y a parfois avantage, dans l'étude de X, de \var{\xout{au}}{}{typewriter strike-out}{ega25_25_055-v0003} remplacer f par $u_n f$ \var{\xout{\unsure{f}}}{}{typewriter strike-out}{ega25_25_055-v0004}
pour éviter un comportement trop particulier et parfois \var{\xout{général}}{}{typewriter strike-out}{ega25_25_055-v0005} gênant .
de f à certains égards. Un exemple d'une telle particularité est celle
\var{\xout{ex}}{}{typewriter strike-out}{ega25_25_055-v0006} signalée (sic) dans 8.12. b), \var{\xout{\unsure{lexmorp}}}{}{typewriter strike-out}{ega25_25_055-v0007} où $\underline{Y}^{\mathrm{sing}} \longrightarrow \overset{\hand{\vee}}{X}$ a une image de
dimension r-1 mais \var{\xout{estxunxmorphismexinsépa}}{}{typewriter strike-out}{ega25_25_055-v0008} donne lieu à une extension
de corps inséparable. Un autre cas est celui qui se présente dans le cas
d'une quadrique \var{\deleted{de dimension impaire \var{\xout{dans}}{}{typewriter strike-out}{ega25_25_055-v0010} 2m lorsque r=2m+1,}}{}{deletion}{ega25_25_055-v0009}\var{\inserted{dans $\mathbf{P}^3$}}{dans $\mathbf{P}^3$}{insertion}{ega25_25_055-v0011} savoir que
toutes les sections hyperplanes singulières sont géométriquement
réductibles (alors que X est cependant géométriquement irréductible).

\uline{Proposition} 9.2. On suppose S=Spec(k), X lisse sur k, et $f\colon X\to P$ non
ramifié. Soit $n \geq 2$, et considérons $f_n = u_n f$ . Alors $f_n : X \to P(n)$
satisfait aux condition équivalentes de 8.8. , en particulier
pour \hand{$\xi \in$} $\overset{\hand{\vee}}{P}(n)$ dans le complémentaire d'un ensemble de codimension $\geq 2$,
la section hyperp\var{\xout{b}}{}{typewriter strike-out}{ega25_25_055-v0012}lane correspondante $Y_{\hand{\xi}}$ est lisse \var{\xout{sur}}{}{typewriter strike-out}{ega25_25_055-v0013} ou \var{\inserted{n'}}{n'}{insertion}{ega25_25_055-v0014}admet qu'un
\var{\inserted[typed]{nombre fini de points non lisses ,}}{nombre fini de points non lisses ,}{insertion}{ega25_25_055-v0015}
\var{\xout{plusxunxseulxpointxnonxlissex}}{}{typewriter strike-out}{ega25_25_055-v0016}, lesquels sont géométriquement singuliers
ordinaires\var{\xout{etxrationnelsxsurxk(xx)x}}{}{typewriter strike-out}{ega25_25_055-v0017}. Si f est une immersion, il y a au
plus un tel point singulier, et il est rationnel sur k(\hand{$\xi$}).
\var{\margin{$\|$}}{\placedmargin{$\|$}}{marginal note placed}{ega25_25_055-v0018} NB Il fallait énoncer 8.8. de façon à \var{\xout{\unsure{éviter}}}{}{typewriter strike-out}{ega25_25_055-v0019} ne pas exclure le
cas où f ne serait pas une immersion. La \unsure{vérifiaation|vérification} est essentielle-%
ment triviale sur la condition (ii bis) de 8.8. -Il faudrait sans doute
expliquer dans 9.1. que les sections hyperplanes de X relativement à
\var{\margin{55}}{\placedmargin{55}}{marginal note placed}{ega25_25_055-v0020} $f_n$ ne sont autres que les``sections''de X par les hypersurfaces de degré
